\documentclass[11pt,a4paper]{article}

% --------------------
% PACKAGES
% --------------------
\usepackage[utf8]{inputenc}
\usepackage[T1]{fontenc}
\usepackage{lmodern}
\usepackage[english]{babel}
\usepackage{geometry}
\usepackage{hyperref}
\usepackage{graphicx}
\usepackage{listings}
\usepackage{xcolor}
\usepackage{fancyhdr}
\usepackage{enumitem}
\usepackage{titlesec}
\usepackage{etoolbox}  % For BeforeBeginEnvironment and AfterEndEnvironment

% --------------------
% PAGE SETTINGS
% --------------------
\geometry{margin=1in}
\pagestyle{fancy}
\fancyhf{}
\fancyhead[L]{\leftmark}
\fancyhead[R]{\thepage}
\setlength{\parskip}{0.5em}
\raggedbottom  % Allow pages to have flexible height to avoid breaking code blocks

% --------------------
% COLORS AND CODE BLOCKS
% --------------------
\definecolor{codegray}{gray}{1}
\definecolor{codeblue}{rgb}{0.2,0.2,0.6}

\lstset{
	backgroundcolor=\color{codegray},
	basicstyle=\ttfamily\small,
	keywordstyle=\color{codeblue}\bfseries,
	commentstyle=\color{gray}\itshape,
	stringstyle=\color{orange},
	numbers=left,
	numberstyle=\tiny,
	stepnumber=1,
	numbersep=10pt,
	showstringspaces=false,
	breaklines=true,
	frame=single,
	captionpos=b,
	language=C
}

% Prevent page breaks within lstlisting environments by wrapping them in minipages
\BeforeBeginEnvironment{lstlisting}{\par\nobreak\minipage{\linewidth}}
\AfterEndEnvironment{lstlisting}{\endminipage}

% --------------------
% TITLE
% --------------------
\title{
	\textbf{\Huge The Serotonin Operating System Manual} \\
	\large{Volume 2: C/C++ Usermode Programming} \\
	\vspace{0.5cm}
	\vspace{0.3cm}
	Version: 0.4.0 \\
	\vspace{0.5cm}
	\textit{} \\
	\date{\today}
}
\author{}

\setcounter{secnumdepth}{4}
\titleformat{\paragraph}
{\normalfont\normalsize\bfseries}{\theparagraph}{1em}{}
\titlespacing*{\paragraph}
{0pt}{3.25ex plus 1ex minus .2ex}{1.5ex plus .2ex}


% --------------------
% DOCUMENT
% --------------------
\begin{document}

	\maketitle
	\tableofcontents
	\newpage

	\section{Getting Started}
	Serotonin is a Unix-like operating system. If you already have experience programming for other Unix-like operating systems (Linux, macOS), especially in C or C++, then you will probably find it easy to program for Serotonin. There does exist non-POSIX system call (lib5ht), but they are designed to be easier to interface with compared to typical Unix-like operating systems.
	Overall, if you just want to have some fun programming on an easy to learn yet powerful operating system, Serotonin is well suited for it.

	\subsection{Tutorial: Hello World}
	Prerequisites:
	\begin{itemize}
		\item A full clone of the Serotonin git repository with all tools built $\texttt{OR}$
		\item A copy of binaries, including all development tools.
	\end{itemize}
	The cross-compiler is especially important, as it is required for compiling programs for Serotonin on other operating systems. It is assumed that you know how to compile a program for Serotonin. If not, please read \texttt{Volume 1: 3.1/3.2} for more guidance.
	\begin{lstlisting}[language=C]
#include <stdio.h>
int main(void) {
    printf("Hello, world!\n");
    return 0;
}
	\end{lstlisting}
	You will notice that this is basically the same for almost all Unix-like operating systems, as Serotonin provides the Newlib C Standard Library.

	\subsection{Example: Process Management}
	\begin{lstlisting}[language=C]
#include <stdio.h>
#include <unistd.h>
int main(int argc, char **argv, char **envp) {
    printf("I am the parent.\n");
    pid_t fork_result = fork();
    if (fork_result < 0) {
        // fork failed
        return 1;
    } else if (fork_result == 0) {
    	printf("I am the child.\n");
    	execve("/bin/sh", {"/bin/sh"}, envp);
    } else if (fork_result > 0) {
        printf("child pid: %d\n", fork_result);
        int status;
        waitpid(fork_result, &status, 0);
    }
    return 0;
}
	\end{lstlisting}
	\subsection{Example: Error Handling}
	\subsection{Example: Inter Process Communication}
	\subsection{Example: SSE2 Acceleration}
	\subsection{Example: Interacting With Devices}
	\subsection{Example: Drawing GUIs}

	\section{System Call Manual}
	\subsection{exit}
    \begin{itemize}
		\item Kernel version: 0.1.0+
		\item Usage: $\texttt{\_exit(int status)$}$
	\end{itemize}

	Terminates the process. This function does not return. This is called by the Serotonin runtime when returning from \texttt{main()}.

	\subsection{write}
    \begin{itemize}
		\item Kernel version: 0.1.0+
		\item Usage: $\texttt{write(int fd, const void* buf, size\_t count)$}$
	\end{itemize}
	Write to a file descriptor. Returns the number of bytes written, or -1 indicating an error.

	\subsection{read}
    \begin{itemize}
		\item Kernel version: 0.1.0+
		\item Usage: $\texttt{read(int fd, void* buf, size\_t count)$}$
	\end{itemize}
	Read from a file descriptor. Returns the number of bytes read, or -1 indicating an error.

	\subsection{execve}
    \begin{itemize}
		\item Kernel version: 0.2.0+
		\item Usage: $\texttt{execve(const char *name, char *const argv[], char *const envp[)$}$
	\end{itemize}
	Replaces the current program image with a new program. This function only returns on error.

	\section{Desktop Enviorment Programming}
\end{document}
